\section{Related Work}

\subsection{Generative Explainable Recommendation and Degeneration}

Explainable recommendation studies models that pair recommendation decisions
with textual or structured rationales \cite{explainable_rec_survey}. Neural
and language-model-based systems such as PETER, PEPLER, ReXPlug, ERRA, and
AdaReX demonstrate that explanations can be generated jointly with
recommendation signals \cite{peter,pepler,rexplug,erra,adarex}. These models
also highlight the central risk for this paper: fluent explanations are not
automatically faithful explanations. When a generator overuses common review
phrases or domain templates, explanation quality degenerates even if the text
looks plausible.

\subsection{Cross-Domain Explainable Recommendation}

Cross-domain recommendation transfers evidence from a source domain to a
target domain \cite{crossdomain_survey}. Cross-domain explainable
recommendation adds a second transfer problem: the explanation must preserve
recommendation content while adapting to target-domain language. This makes
content and style inseparable if the model treats text as one coarse signal.
ODCR addresses this by representing content evidence and style state as
separate control variables rather than relying on a single domain-level
textual attribute.

\subsection{Causal and Counterfactual Recommendation}

D4C is the closest causal foundation for ODCR \cite{d4c}. It diagnoses
cross-domain explanation degeneration through hidden textual attributes \(W\)
and uses domain counterfactual reasoning to reduce shortcut dependence. ODCR
inherits this causal view but emphasizes the dynamic feedback part of the
diagnosis: \(W^{(t)}\) influences explanations, and explanations can feed back
into future textual attributes. ODCR refines this process into stable content
evidence \(C\), dynamic style \(S^{(t)}\), and reliability-gated
counterfactual influence \(R_{cf}\).

\subsection{Controlled Generation, Faithfulness, and Consistency}

Controlled explanation generation requires more than producing text in the
right domain style. The generated explanation should remain tied to the
content evidence that supports the recommendation. ODCR therefore positions
controlled generation and faithfulness alignment as output-side parts of the
method: \(C\) controls what is said, \(S^{(t)}\) controls how it is expressed,
and \(R_{cf}\) controls which counterfactual signals are trusted during
learning. Modern external systems such as CIER, MAPLE, ELIXIR, Prism, and XRec
remain comparison targets once adapted to the same data and evaluation setting;
the current manuscript does not claim results against them.
